\documentclass[twocolumn,a4j,dvipdfmx]{jsarticle_customized1}

\usepackage{amssymb,amsmath,bm}
\usepackage{graphicx,tabularx}
\usepackage{float}
\usepackage{times,type1cm}
\usepackage{url}
\usepackage{color}
\setlength{\abovecaptionskip}{0pt}
\setlength{\belowcaptionskip}{0pt}
\setlength{\intextsep}{5pt}

\title{介護マッチングサービスにおけるマッチング精度向上の提案と実装}
\author{星子 祐哉}
\date{}
\labname{情報メディア環境学研究室}

\begin{document}
\twocolumn[\maketitle\vspace{0.0zw}]
\setlength{\baselineskip}{12.6pt}
\setlength{\columnsep}{8pt}
\setlength{\parindent}{1zw}

 \section{研究背景と課題}
介護業界では人材不足が続いており、2040年問題（高齢者人口の急増と生産年齢人口の減少が同時に進む局面）により、介護需要の拡大と人材不足の深刻化が予想される。中小事業者では1人採用するのに広告・面接・研修などの費用がかさみ、本来なら2〜3人を雇える水準の費用になるため、早期離職を防ぐことが経営の死活問題になる。とりわけ北海道のような高齢化率33.5\%超・人口偏在地域では、札幌圏集中と交通インフラ不足により応募が偏り、地方事業所は近隣でしか人材を確保できない。加えて、積雪や凍結の影響もあり、通勤条件が応募に与える影響は大きい。

こうした課題に対し、マッチングアプリは求人とワーカーを効率的につなぐ手段として普及しつつある。しかし、従来の求人検索は勤務地・資格・給与といった定型条件の一致に偏り、職種や業務内容に対する意味的な親和性が評価されにくい。介護分野では業務の幅が広く、同じ職種名でも現場ごとの差が大きいため、条件が合っていても「実際の業務像」が異なり、応募後のギャップが生じやすい。

さらに、ワーカーには「入浴介助を避けたい」「夜勤を希望しない」といった文脈的な志向がある。ところが、現行の検索はこれらを捉えにくく、条件一致でも実務ギャップが生じ離職につながる。特定施設での経験やスタッフ構成といった暗黙的な希望も多いにもかかわらず、求人票の雰囲気や業務ニュアンスなど非構造化情報は十分に活用されていない。

そこで本研究では、求人情報の意味的な近さと応募履歴・嗜好の変化・回避傾向を統合し、表面的な条件一致に依存しない推薦を目指す。具体的な類似度計算や重み付けの設計は次章で述べる。


\section{提案手法}
提案手法では、求人タイトル・業務内容・介護項目・施設名・勤務条件など複数のテキスト属性を自然言語処理モデルにより数値ベクトルに変換し、重み付きコサイン類似度により職種的・スキル的親和性をスコア化する。各テキスト属性に重みを付与して類似度を統合することで、単一キーワード一致では把握しきれない文脈的な近さを捉える。特に業務内容と条件系の文脈が、応募後の満足度に影響する点を重視した。

また、過去応募履歴を嗜好ベクトルとして平均化し、時間的減衰（指数・線形・ウィンドウ・セッション順）を導入して直近履歴を強調する。時間重みの自動選別では、嗜好の変化が大きなワーカーには半減期7日の指数減衰を、安定志向のワーカーには緩やかな線形重みを選び、新旧履歴のバランスを取る。

\section{実験と評価}
実験環境はPythonで作成し、FastAPIを用いた\cite{fastapi}。推薦結果はTop-1/Top-3/Top-5を基本指標とし、実運用で「おすすめを1件提示する」要件を踏まえて複合スコアを併記する。複合スコア $S_{\text{composite}}$ はTop-1/3/5を重み付きで統合した指標で、以下の式で定義する。
\[
S_{\text{composite}}=0.5\,\text{Top-1}+0.3\,\text{Top-3}+0.2\,\text{Top-5}
\]
予備実験を行い、CLS埋め込み・mean pooling・Sentence-BERTについて分離度を比較した結果、最も分離度の高かった日本語Sentence-BERTを採用した。分離度は「類似職種の平均スコア−非類似職種の平均スコア」で定義され、値が大きいほど類似/非類似をより明確に区別できる。

 \subsection{実験1：均等重みとTop-$k$のベースライン}
本実験1では7項目（職種名、業務内容、介護項目、施設、勤務条件など）をすべて同じ重要度として扱い、各項目の類似度を均等に足し合わせて求人の順位付けを行った。18〜20名の応募履歴を15〜30件で学習した。データは候補提示・応募のペアを1セッション単位として扱い、先頭の10件を固定テストとした。均等重みはベースラインとして機能する。
 \begin{table}[H]
 \caption{実験1：学習件数別Top-$k$精度}
 \label{tab:exp1-topk}
 \centering
 \resizebox{\linewidth}{!}{%
 \begin{tabular}{lcccc}
 \hline
 学習件数 & Top-1 & Top-3 & Top-5 & 複合スコア \\
 \hline
 15 & 0.517 & 0.600 & 0.678 & 0.574 \\
 20 & 0.478 & 0.628 & 0.694 & 0.566 \\
 30 & 0.309 & 0.418 & 0.527 & 0.386 \\
 \hline
 \end{tabular}%
 }
 \end{table}
15〜20件が最も安定し、30件ではTop-1が0.365まで低下したため、一定件数以降の履歴を抑えて直近を重視する設計が望ましい。

 \subsection{実験2：項目別/時間重みの探索}
実験1では学習件数が増えると精度が低下する傾向が見られたため、実験2では履歴の新しさを反映する時間重み（指数・線形・ウィンドウ・セッション順）を変化させ、嗜好ドリフトへの耐性を評価した。あわせて項目別重みの違いも比較し、影響度を確認した。指数減衰の例として、時刻差 $\Delta t$ に対する重み $w_t$ を次式で定義する。\[
w_t=\exp(-\lambda\, \Delta t)
\]
平均Top-$k$値に加え複合スコアも用いて1件推薦の安定度を追跡した。Top-1だけでなく候補全体の安定性を評価するためである。

 \begin{table}[H]
 \caption{実験2：代表的時間重みによる複合スコア}
 \label{tab:exp2-timeweights}
 \centering
 \resizebox{\linewidth}{!}{%
 \begin{tabular}{lcccc}
 \hline
 時間重み & 複合スコア & Top-1 & Top-3 & Top-5 \\
 \hline
 直近7日指数 & 0.634 & 0.556 & 0.694 & 0.739 \\
 直近30日線形 & 0.614 & 0.533 & 0.672 & 0.728 \\
 時間重みなし & 0.566 & 0.478 & 0.628 & 0.694 \\
 \hline
 \end{tabular}%
 }
 \end{table}
直近7日の指数減衰が最良スコアを記録し、直近30日線形も好成績であることから、古い履歴を弱める設計全般が精度向上に寄与する。また、業務内容・介護項目・条件系を同時に強化した項目重みが平均精度の安定化に貢献した。短期志向のワーカーほど効果が大きい傾向が見られた。

\subsection{実験3：時間志向自動選択の実装と結果}
実験3では、履歴の変化量（drift）に応じて時間重みを自動で切り替える仕組みを実装し、時間重みなしおよび固定の直近30日線形減衰と比較した。drift は直近履歴ベクトルと過去履歴ベクトルのコサイン類似度から算出し、$ \mathrm{drift}=1-\cos(\bm{v}_{\mathrm{recent}},\bm{v}_{\mathrm{past}})$ と定義する。drift が 0.12 以上なら短期志向（直近重視の指数減衰）に、0.12 未満なら長期志向（期間制限型）に判定する。18名・学習20件・特徴量配分均等の条件で評価した結果を表\ref{tab:time-auto-results}に示す。時間重み自動判別は時間減衰なしより改善した一方で、直近30日線形減衰のほうが高い値となった。
\begin{table}[H]
\caption{時間重み自動判別の実装結果（18名・学習20件）}
\label{tab:time-auto-results}
\centering
\resizebox{\linewidth}{!}{%
\begin{tabular}{lcccc}
\hline
時間重み & Top-1 & Top-3 & Top-5 & 複合スコア \\
\hline
時間減衰なし & 0.478 & 0.628 & 0.694 & 0.566 \\
時間重み自動判別 & 0.494 & 0.650 & 0.706 & 0.577 \\
直近30日線形減衰 & 0.533 & 0.672 & 0.728 & 0.614 \\
\hline
\end{tabular}%
}
\vspace{-0.5em}
\end{table}


\subsection{実装と運用}
評価システムはGUI/CLIのどちらからも同一評価器を呼び出せる構成で、Jinja2ベースのテンプレートをFastAPIで提供する\cite{jinja2,fastapi}。非同期タスクで複数重みの並列評価を行い、結果はJSON/HTML/CSVに記録する。進捗ログや部分結果の保存により再現性を担保した。求人埋め込みは事前計算とキャッシュで保持し、Apple Silicon向けMPS\cite{mps}やスレッド最適化で再実行時間を短縮した。

\subsubsection*{評価パイプライン}
データ読み込み→前処理→埋め込み→スコアリング→推薦生成→出力保存の流れで処理し、Top-$k$・複合スコア・学習曲線・推定スコアの分散などを記録する。

\section{ユーザースタディ}
インタビュー2件（副業で月4〜5回利用する介護経験7年目ワーカーと、実務者研修経由で単発バイトを始めた利用者）から、ワーカーにはリピート志向が強く、一度良い施設には繰り返し応募する傾向が確認された。経験が長いほど「慣れた施設で安定的に働きたい」という志向が強い。また、往復交通費が500円を超えると手取りが悪化するため、駐車場や送迎などの非定型条件が意思決定に直結する。入浴介助等の特定業務は時給よりも強い回避条件として働き、情報収集は「近隣→業務内容→口コミ」の順で行われる傾向が見られた。

さらに、回避項目と総額志向（1回あたり約1万円）の関係や初回訪問時の情報非対称性が示唆され、回避業務が並ぶと応募に至りにくいという実務的課題が明らかになった。これらの知見はペナルティ係数の設計に直結し、避けたい項目を可視化することで判断負荷を下げ、継続利用につながると考えられる。説明性の高い提示が応募の質に影響する点も示唆された。

\section{追加実験：回避ペナルティの効果}
実験目的は、ユーザースタディで判明した「特定の業務（入浴介助など）を避けたい」というワーカーの選好を推薦アルゴリズムに反映させ、精度の向上を検証することである。
手法として、候補求人と選択求人のタグを集計し、選択率30\%未満を回避項目と判定したうえで、回避スコアに基づきペナルティ重み0.5を乗算して再ランキングする。
使用データは、実際のワーカー18名分の行動履歴データである。

表\ref{tab:penalty-simple}に示す通り、回避ペナルティを適用した条件で精度が改善した。これにより、回避傾向をスコアに直接反映することが推薦精度の向上に寄与することが示唆された。

\begin{table}[H]
\caption{回避ペナルティの効果（学習20件・18名・時間重みなし）}
\label{tab:penalty-simple}
\centering
\resizebox{\linewidth}{!}{%
\begin{tabular}{lcccc}
\hline
条件 & Top-1 & Top-3 & Top-5 & 複合スコア \\
\hline
ペナルティ無し & 0.478 & 0.628 & 0.694 & 0.566 \\
ペナルティ有り & 0.583 & 0.744 & 0.806 & 0.669 \\
差分 & +0.106 & +0.117 & +0.111 & +0.103 \\
\hline
\end{tabular}%
}
\vspace{-0.5em}
\end{table}

回避ペナルティの付与方法は、求人に含まれる回避項目の割合と強さを用いて次式で計算する。回避項目の割合を $r$、回避項目の最大回避スコアを $s_{\max}$、重みを $w=0.5$ とすると、
\[
\text{penalty\_factor}=(1-wr)\times(1-ws_{\max})
\]
と定義する。元の推薦スコアに \(\text{penalty\_factor}\) を乗じて再ランキングすることで、回避項目を多く含む求人や回避が強い項目を含む求人ほど順位が下がる。

\section{まとめと展望}
本研究では、意味的類似性と行動履歴を統合したマッチングを設計し、学習データ量や時間重みの与え方で精度が変動することを示した。特に、履歴の変化を考慮した時間重みや回避業務ペナルティの導入により、Top-$k$精度の改善が確認された。一方で、データ規模や入力品質のばらつき、回避傾向の推定精度などは課題として残るため、実運用に近い条件での検証と改善サイクルが必要である。そこで、これらの課題を踏まえた今後の展望を以下に整理する。
\paragraph{ユーザースタディの拡大}
ユーザースタディの件数を増やし、より多角的な定性的視座を得る。今回は2件のみだったが、それでも多くの知見を得られたため、今後は属性や就業形態の異なる参加者を広く募り、回避項目や意思決定要因の多様性を整理する。

\paragraph{リアルタイムフィードバックの実装と評価}
クリック回数や閲覧時間など異なる角度の指標を取り入れる実装を行い、それらを加えた上で重み付けを再設計し、効果を検証する。暗黙的フィードバックを用いることで、応募に至る前段階の関心や回避傾向も反映できる可能性がある。

\paragraph{実システム実装によるマッチング数の変移の検証}
実際に実装したうえで、ユーザーの行動変化とマッチング数の推移を確認する。推薦の提示方法や説明性の違いが応募率や継続利用に与える影響も併せて確認し、実運用に近い条件での有効性を評価する。

\vspace{-0.8em}
{\scriptsize
\setlength{\itemsep}{0pt}
\setlength{\parsep}{0pt}
\renewcommand{\refname}{\small 参考文献}
\begin{thebibliography}{9}
\bibitem{fastapi} FastAPI Documentation, \url{https://fastapi.tiangolo.com/}
\bibitem{jinja2} Jinja Documentation, \url{https://jinja.palletsprojects.com/}
\bibitem{mps} Apple Developer Documentation, Metal Performance Shaders, \url{https://developer.apple.com/documentation/metalperformanceshaders}
\end{thebibliography}
}

\vspace{-0.6em}
\section*{\small 著者の研究業績}
{\scriptsize
\setlength{\itemsep}{0pt}
\setlength{\parsep}{0pt}
\begin{enumerate}
\item 星子 祐哉, 青木 直史, 土橋 宜典 “VGG16 の転移学習を利用した病的音声の検出”，電気・情報関係学会北海道支部連合大会，オンライン開催，November 5-6, 2022．
\item 介護マッチングサービスにおけるマッチング精度向上の提案と実装（星子 祐哉, 土橋 宜典），令和 7 年度 電気・情報関係学会北海道支部連合大会（2025/11/2 - 3）
\end{enumerate}
}
\end{document}
